\documentclass[12pt]{article}

% هوامش
\usepackage{geometry}
\geometry{margin=0.8in}

% ألوان، رؤوس/تذييل، تباعد، قوائم، أعمدة
\usepackage{xcolor}
\usepackage{fancyhdr}
\usepackage{setspace}
\usepackage{enumitem}
\usepackage{multicol}
\columnsep=1.0cm

% دعم يونيكود ولغات
\usepackage{fontspec}
\usepackage{polyglossia}

% الاتجاه العام عربي (بدل إلى english إن رغبت)
% اختر نوع الأرقام: western (0-9) أو eastern (٠-٩)
\setmainlanguage[numerals=western]{arabic}
\setotherlanguage{english}

% الخطوط (بدّل إذا لزم)
\newfontfamily\arabicfont[Script=Arabic,Scale=1.05]{Amiri}
\newfontfamily\englishfont{Times New Roman}

% ترويسة/تذييل
\pagestyle{fancy}
\fancyhf{}
\renewcommand{\headrulewidth}{0pt}
\cfoot{\textenglish{Page } \thepage}

% ماكرو خيار
\newcommand{\choice}[3]{%
  \noindent\textbf{#1}.~#2\hfill#3\par\vspace{0.25em}
}

% إعدادات فقرات
\setlength{\parskip}{0.4em}
\setlength{\parindent}{0pt}

\begin{document}

% رأس المستند
\begin{flushright}
{\Large \textbf{{\textenglish{Algorithms -1-}}}}\\[0.25em]
{\normalsize \textenglish{Model: }{\textenglish{D}}}\\[0.1em]
{\normalsize \textenglish{Date: }{\textenglish{September 4, 2025}}}
\end{flushright}

\vspace{0.6cm}

% معلومات إضافية
\noindent
\begin{minipage}[t]{0.65\textwidth}
  \raggedright
  \textbf{{\textenglish{University of Aleppo}}}\\[0.3em]
\end{minipage}\hfill
\begin{minipage}[t]{0.3\textwidth}
  \raggedleft
  \textenglish{\textbf{Duration: }}{\textenglish{90 minuts}}
\end{minipage}

\vspace{0.9cm}

\section*{اختر الإجابة الصحيحة لكل سؤال}
\setstretch{1.2}

\begin{multicols}{2}
\noindent \textbf{\textenglish{1.}} {\textenglish{Hello!}}

\noindent\choice{A}{\underline{{\textenglish{1}}}}{}
\noindent\choice{B}{{\textenglish{2}}}{}
\noindent\choice{C}{{\textenglish{3}}}{}
\noindent\choice{D}{{\textenglish{4}}}{}
\hfill \textbf{[2 points]}\vspace{0.5em}

\noindent \textbf{\textenglish{2.}} {\textenglish{Test}}

\noindent\choice{A}{{\textenglish{13}}}{}
\noindent\choice{B}{\underline{{\textenglish{14}}}}{}
\noindent\choice{C}{{\textenglish{15}}}{}
\noindent\choice{D}{{\textenglish{16}}}{}
\noindent\choice{E}{{\textenglish{17}}}{}
\hfill \textbf{[3 points]}\vspace{0.5em}

\noindent \textbf{\textenglish{3.}} {\textenglish{\textbackslash{}(x=3+6 \textbackslash{}) what is the value of x?}}

\noindent\choice{A}{{\textenglish{7}}}{}
\noindent\choice{B}{\underline{{\textenglish{9}}}}{}
\noindent\choice{C}{{\textenglish{8}}}{}
\noindent\choice{D}{{\textenglish{10}}}{}
\noindent\choice{E}{{\textenglish{0}}}{}
\hfill \textbf{[1 points]}\vspace{0.5em}

\noindent \textbf{\textenglish{4.}} {\textenglish{Test3}}

\noindent\choice{A}{{\textenglish{1}}}{}
\noindent\choice{B}{{\textenglish{2}}}{}
\noindent\choice{C}{\underline{{\textenglish{3}}}}{}
\noindent\choice{D}{{\textenglish{4}}}{}
\noindent\choice{E}{{\textenglish{5}}}{}
\hfill \textbf{[1 points]}\vspace{0.5em}

\noindent \textbf{\textenglish{5.}} {\textenglish{Please calculate \textbackslash{}(\textbackslash{}int f(x)dx\textbackslash{})}}

\noindent\choice{A}{{\textenglish{13}}}{}
\noindent\choice{B}{\underline{{\textenglish{14}}}}{}
\noindent\choice{C}{{\textenglish{15}}}{}
\noindent\choice{D}{{\textenglish{16}}}{}
\hfill \textbf{[2 points]}\vspace{0.5em}

\noindent \textbf{\textenglish{6.}} {\textarabic{احسب مشتق التابع \textbackslash{}( f(x) = x\^{}2\textbackslash{}) وأوجد نهايته عند اللانهاية \textbackslash{}(\textbackslash{}lim\_{x \textbackslash{}to \textbackslash{}infty}f(x)\textbackslash{})}}

\noindent\choice{A}{{\textenglish{14}}}{}
\noindent\choice{B}{{\textenglish{15}}}{}
\noindent\choice{C}{\underline{{\textenglish{16}}}}{}
\noindent\choice{D}{{\textenglish{17}}}{}
\noindent\choice{E}{{\textenglish{17}}}{}
\hfill \textbf{[1 points]}\vspace{0.5em}

\noindent \textbf{\textenglish{7.}} {\textenglish{test4}}

\noindent\choice{A}{\underline{{\textenglish{1}}}}{}
\noindent\choice{B}{{\textenglish{2}}}{}
\noindent\choice{C}{{\textenglish{3}}}{}
\noindent\choice{D}{{\textenglish{4}}}{}
\noindent\choice{E}{{\textenglish{5}}}{}
\hfill \textbf{[1 points]}\vspace{0.5em}

\noindent \textbf{\textenglish{8.}} {\textenglish{Find the inverse of the following matrix: \textbackslash{}(\textbackslash{}begin{bmatrix} 1 \& 3 \& 4 \textbackslash{}\textbackslash{} 5 \& 2 \& 1 \textbackslash{}\textbackslash{} 6 \& 8 \& 1 \textbackslash{}end{bmatrix}\textbackslash{})}}

\noindent\choice{A}{\underline{{\textenglish{There is no inverse for this matrix}}}}{}
\noindent\choice{B}{{\textenglish{\textbackslash{}[\textbackslash{}begin{pmatrix} 1 \& 2 \& 3 \textbackslash{}\textbackslash{} 4 \& 5 \& 6 \textbackslash{}\textbackslash{} 7 \& 8 \& 9  \textbackslash{}end{pmatrix}\textbackslash{}]}}}{}
\noindent\choice{C}{{\textenglish{0}}}{}
\noindent\choice{D}{{\textenglish{0}}}{}
\hfill \textbf{[4 points]}\vspace{0.5em}

\noindent \textbf{\textenglish{9.}} {\textenglish{Solve for \textbackslash{}(x \textbackslash{}) in the equation \textbackslash{}(2x+6 = 0 \textbackslash{})}}

\noindent\choice{A}{{\textenglish{-1}}}{}
\noindent\choice{B}{{\textenglish{-2}}}{}
\noindent\choice{C}{\underline{{\textenglish{-3}}}}{}
\noindent\choice{D}{{\textenglish{3}}}{}
\hfill \textbf{[1 points]}\vspace{0.5em}

\noindent \textbf{\textenglish{10.}} {\textenglish{Test2}}

\noindent\choice{A}{\underline{{\textenglish{1}}}}{}
\noindent\choice{B}{{\textenglish{2}}}{}
\noindent\choice{C}{{\textenglish{3}}}{}
\noindent\choice{D}{{\textenglish{4}}}{}
\noindent\choice{E}{{\textenglish{5}}}{}
\hfill \textbf{[3 points]}\vspace{0.5em}


\end{multicols}

\end{document}
